\documentclass[11pt,a4paper]{article}
\usepackage[utf8]{inputenc}
\usepackage[italian]{babel}
\usepackage[T1]{fontenc}
\usepackage{geometry}
\usepackage{xcolor}
\usepackage{tabularx}
\usepackage{booktabs}
\usepackage{listings}
\usepackage{tikz}
\usepackage{pgf-umlsd}
\usepackage{float}
\usepackage{hyperref}
\usepackage{amsmath}
\usepackage{enumitem}
\usepackage{tcolorbox}

\geometry{margin=2.5cm}

% ============================================
% COLORI E STILI
% ============================================

\definecolor{codebg}{RGB}{40, 42, 54}
\definecolor{codefg}{RGB}{248, 248, 242}
\definecolor{codecomment}{RGB}{98, 114, 164}
\definecolor{codekeyword}{RGB}{139, 233, 253}
\definecolor{codestring}{RGB}{255, 184, 108}
\definecolor{codefunction}{RGB}{80, 250, 123}
\definecolor{accent}{RGB}{139, 92, 246}

\lstdefinestyle{typescript}{
    language=TypeScript,
    backgroundcolor=\color{codebg},
    basicstyle=\ttfamily\small\color{codefg},
    commentstyle=\color{codecomment}\itshape,
    keywordstyle=\color{codekeyword}\bfseries,
    stringstyle=\color{codestring},
    functionstyle=\color{codefunction},
    numberstyle=\tiny\color{codecomment},
    breakatwhitespace=false,
    breaklines=true,
    captionpos=b,
    keepspaces=true,
    numbers=left,
    numbersep=5pt,
    showspaces=false,
    showstringspaces=false,
    showtabs=false,
    tabsize=2,
    frame=single,
    rulecolor=\color{accent},
    frameround=tttt,
}

\lstdefinestyle{javascript}{
    language=JavaScript,
    backgroundcolor=\color{codebg},
    basicstyle=\ttfamily\small\color{codefg},
    commentstyle=\color{codecomment}\itshape,
    keywordstyle=\color{codekeyword}\bfseries,
    stringstyle=\color{codestring},
    numberstyle=\tiny\color{codecomment},
    breakatwhitespace=false,
    breaklines=true,
    captionpos=b,
    keepspaces=true,
    numbers=left,
    numbersep=5pt,
    showspaces=false,
    showstringspaces=false,
    showtabs=false,
    tabsize=2,
    frame=single,
    rulecolor=\color{accent},
    frameround=tttt,
}

% ============================================
% METADATA
% ============================================

\title{
    \textbf{Documentazione Tecnica: FlashcardItem Component} \\
    \large Sistema STUDY - Modulo Flashcard
}
\author{Sistema ExtraTracker}
\date{\today}

\hypersetup{
    colorlinks=true,
    linkcolor=accent,
    filecolor=accent,
    urlcolor=accent,
    citecolor=accent,
    pdftitle={FlashcardItem Component Documentation},
    pdfauthor={ExtraTracker System},
}

% ============================================
% DOCUMENTO
% ============================================

\begin{document}

\maketitle

\begin{abstract}
Questo documento fornisce un'analisi tecnica completa del componente \texttt{FlashcardItem}, parte del sistema STUDY per la gestione di flashcard. Il componente implementa un'interfaccia dual-mode (visualizzazione e modifica) con supporto per navigazione sorgente, validazione input e gestione stato ottimizzata. La documentazione include analisi architetturale, specifiche delle interfacce, diagrammi di flusso e pattern di interazione.
\end{abstract}

\tableofcontents
\newpage

% ============================================
% SEZIONE 1: OVERVIEW
% ============================================

\section{Overview del Componente}

\subsection{Responsabilità Principali}

Il componente \texttt{FlashcardItem} rappresenta l'unità atomica di visualizzazione e modifica di una singola flashcard all'interno del sistema STUDY. Le responsabilità principali includono:

\begin{itemize}
    \item \textbf{Visualizzazione}: Rendering della card in modalità read-only con domanda (Q) e risposta (A)
    \item \textbf{Modifica}: Interfaccia di editing inline con validazione real-time
    \item \textbf{Navigazione Sorgente}: Integrazione con sistema PDF per jump-to-source
    \item \textbf{Gestione Stato}: Coordinamento tra stato locale e stato del parent
    \item \textbf{Validazione}: Controllo integrità dati prima del salvataggio
\end{itemize}

\subsection{Modello Mentale: Vista Utente vs. Sistema}

\begin{tcolorbox}[colback=blue!5!white,colframe=blue!75!black,title=Modello Mentale]
\textbf{Vista Utente:}
L'utente vede una card con domanda e risposta. Può cliccare per modificare, eliminare o navigare alla fonte nel PDF. Durante la modifica, vede un form con due campi di testo e pulsanti Salva/Annulla.

\textbf{Sistema Interno:}
Il componente gestisce due modalità distinte (ViewMode/EditMode) con state machine. Valida input in tempo reale, sincronizza con il parent tramite callbacks, e gestisce transizioni animate tra stati. Il sistema traccia modifiche non salvate e previene operazioni concorrenti.
\end{tcolorbox}

\subsection{Architettura Modulare}

Il componente è stato refactorizzato seguendo principi Clean Code in una struttura modulare:

\begin{itemize}
    \item \texttt{FlashcardItem.tsx} - Componente orchestratore principale
    \item \texttt{FlashcardItem.ViewMode.tsx} - Modalità visualizzazione
    \item \texttt{FlashcardItem.EditMode.tsx} - Modalità modifica
    \item \texttt{FlashcardItem.constants.ts} - Costanti centralizzate
    \item \texttt{FlashcardItem.types.ts} - Definizioni TypeScript
    \item \texttt{FlashcardItem.utils.ts} - Funzioni pure di utilità
\end{itemize}

% ============================================
% SEZIONE 2: INTERFACCIA E PROPS
% ============================================

\section{Interfaccia e Props (Input)}

\subsection{Tabella Props Completa}

\begin{table}[H]
\centering
\begin{tabularx}{\textwidth}{|l|X|c|c|}
\hline
\textbf{Prop} & \textbf{Descrizione} & \textbf{Tipo} & \textbf{Req.} \\
\hline
\hline
\texttt{card} & Oggetto Card contenente i dati della flashcard (front, back, id, metadata). Rappresenta l'entità di dominio principale. & \texttt{Card} & \checkmark \\
\hline
\texttt{onUpdate} & Callback asincrono invocato quando l'utente salva modifiche. Riceve \texttt{(cardId: string, front: string, back: string)} e restituisce \texttt{Promise<void>}. Gestito dal parent per persistenza. & \texttt{(cardId: string, front: string, back: string) => Promise<void>} & \checkmark \\
\hline
\texttt{onClick} & Callback opzionale invocato quando l'utente clicca sulla card (solo in view mode). Utile per navigazione o selezione. & \texttt{(card: Card) => void} & \textcolor{gray}{Opz.} \\
\hline
\texttt{onDelete} & Callback opzionale per eliminazione card. Riceve \texttt{cardId: string}. Se assente, il pulsante delete non viene renderizzato. & \texttt{(cardId: string) => void} & \textcolor{gray}{Opz.} \\
\hline
\texttt{initialEditing} & Flag booleano per avviare il componente in modalità editing. Default: \texttt{false}. Utile per card temporanee o creazione inline. & \texttt{boolean} & \textcolor{gray}{Opz.} \\
\hline
\texttt{onCancel} & Callback opzionale invocato quando editing viene annullato. Utile per cleanup di card temporanee. & \texttt{() => void} & \textcolor{gray}{Opz.} \\
\hline
\texttt{onCreate} & Callback opzionale per creazione nuova card (alternativa a \texttt{onUpdate}). Usato quando \texttt{card.id} inizia con \texttt{"temp-"}. & \texttt{(front: string, back: string) => Promise<void>} & \textcolor{gray}{Opz.} \\
\hline
\texttt{onShowSource} & Callback opzionale per navigazione alla fonte PDF. Invocato quando utente clicca pulsante "Show Source". Richiede \texttt{card.sourceMetadata} presente. & \texttt{(card: Card) => void} & \textcolor{gray}{Opz.} \\
\hline
\texttt{isSourceActive} & Flag booleano che indica se la card è attualmente evidenziata come "source active" (navigazione PDF attiva). Modifica styling per feedback visivo. & \texttt{boolean} & \textcolor{gray}{Opz.} \\
\hline
\end{tabularx}
\caption{Tabella completa delle props di FlashcardItem}
\label{tab:props}
\end{table}

\subsection{Flusso di Dati dal Parent}

Il componente riceve dati attraverso un pattern unidirezionale:

\begin{enumerate}
    \item \textbf{Props Immutabili}: Il parent passa \texttt{card} come prop immutabile
    \item \textbf{State Locale}: Il componente crea copie temporanee (\texttt{tempFront}, \texttt{tempBack}) per editing
    \item \textbf{Sincronizzazione}: \texttt{useEffect} sincronizza state locale quando \texttt{card} cambia (solo se non in editing)
    \item \textbf{Callback Upstream}: Modifiche vengono propagate al parent tramite \texttt{onUpdate} o \texttt{onCreate}
\end{enumerate}

% ============================================
% SEZIONE 3: GESTIONE STATO E LOGICA
% ============================================

\section{Gestione dello Stato e Logica Interna}

\subsection{State Management}

Il componente utilizza quattro hook di stato:

\begin{lstlisting}[style=typescript, caption=State Hooks]
const [isEditing, setIsEditing] = useState(initialEditing);
const [tempFront, setTempFront] = useState(card.front);
const [tempBack, setTempBack] = useState(card.back);
const [isSaving, setIsSaving] = useState(false);
\end{lstlisting}

\subsubsection{Analisi Dettagliata degli State}

\begin{description}
    \item[\texttt{isEditing}] \textbf{Tipo}: \texttt{boolean}. \textbf{Scopo}: Flag che determina quale modalità renderizzare (ViewMode vs EditMode). \textbf{Transizioni}: \texttt{false} $\rightarrow$ \texttt{true} tramite \texttt{handleStartEdit()}, \texttt{true} $\rightarrow$ \texttt{false} tramite \texttt{handleSave()} o \texttt{handleCancel()}.
    
    \item[\texttt{tempFront}] \textbf{Tipo}: \texttt{string}. \textbf{Scopo}: Buffer temporaneo per il testo della domanda durante editing. \textbf{Sincronizzazione}: Aggiornato da \texttt{card.front} quando si esce da editing mode.
    
    \item[\texttt{tempBack}] \textbf{Tipo}: \texttt{string}. \textbf{Scopo}: Buffer temporaneo per il testo della risposta durante editing. \textbf{Sincronizzazione}: Aggiornato da \texttt{card.back} quando si esce da editing mode.
    
    \item[\texttt{isSaving}] \textbf{Tipo}: \texttt{boolean}. \textbf{Scopo}: Flag per prevenire operazioni concorrenti durante salvataggio. \textbf{Lifecycle}: \texttt{false} $\rightarrow$ \texttt{true} all'inizio di \texttt{handleSave()}, \texttt{true} $\rightarrow$ \texttt{false} nel \texttt{finally} block.
\end{description}

\subsection{Effects e Sincronizzazione}

\subsubsection{useEffect per Sincronizzazione}

\begin{lstlisting}[style=typescript, caption=Sincronizzazione State]
useEffect(() => {
    if (!isEditing) {
        setTempFront(card.front);
        setTempBack(card.back);
    }
}, [card.front, card.back, isEditing]);
\end{lstlisting}

\textbf{Logica}: L'effect sincronizza i valori temporanei con i valori originali della card \textbf{solo quando non si è in modalità editing}. Questo previene perdita di dati durante editing e garantisce coerenza quando la card viene aggiornata esternamente.

\textbf{Dipendenze}: 
\begin{itemize}
    \item \texttt{card.front}, \texttt{card.back}: Trigger quando il parent aggiorna la card
    \item \texttt{isEditing}: Condizione di guardia per evitare sovrascritture durante editing
\end{itemize}

\subsection{Validazione e Computed Values}

\subsubsection{Funzione di Validazione}

La validazione è implementata tramite funzione pura in \texttt{FlashcardItem.utils.ts}:

\begin{lstlisting}[style=typescript, caption=Validazione Contenuto]
export function validateCardContent(
    tempFront: string,
    tempBack: string,
    originalFront: string,
    originalBack: string,
    isSaving: boolean
): ValidationResult {
    const trimmedFront = tempFront.trim();
    const trimmedBack = tempBack.trim();
    const isDirty = trimmedFront !== originalFront || 
                    trimmedBack !== originalBack;
    const canSave = !isSaving && 
                   trimmedFront.length > 0 && 
                   trimmedBack.length > 0 && 
                   isDirty;
    return { trimmedFront, trimmedBack, isDirty, canSave };
}
\end{lstlisting}

\textbf{Regole di Validazione}:
\begin{enumerate}
    \item \textbf{Trim}: Rimozione spazi iniziali/finali
    \item \textbf{Length Check}: Entrambi i campi devono avere lunghezza $> 0$
    \item \textbf{Dirty Check}: Almeno un campo deve essere modificato rispetto all'originale
    \item \textbf{Concurrency Check}: Salvataggio non in corso
\end{enumerate}

\subsection{Event Handlers e Callbacks}

\subsubsection{handleStartEdit}

\begin{lstlisting}[style=typescript, caption=Handler Inizio Editing]
const handleStartEdit = useCallback(() => {
    setTempFront(card.front);
    setTempBack(card.back);
    setIsEditing(true);
}, [card.front, card.back]);
\end{lstlisting}

\textbf{Comportamento}: 
\begin{enumerate}
    \item Reset dei valori temporanei ai valori originali (previene carry-over di modifiche precedenti)
    \item Transizione a modalità editing
    \item \texttt{useCallback} previene re-render inutili dei children
\end{enumerate}

\subsubsection{handleSave}

\begin{lstlisting}[style=typescript, caption=Handler Salvataggio]
const handleSave = useCallback(async () => {
    if (!validation.canSave) return;
    setIsSaving(true);
    try {
        if (isTemporaryCard(card.id) && onCreate) {
            await onCreate(validation.trimmedFront, validation.trimmedBack);
        } else {
            await onUpdate(card.id, validation.trimmedFront, validation.trimmedBack);
            emitToast.success(TEXT_CONTENT.toast.success.message, {...});
        }
        setIsEditing(false);
    } catch (error) {
        emitToast.error(errorMessage, {...});
    } finally {
        setIsSaving(false);
    }
}, [card.id, validation, onUpdate, onCreate]);
\end{lstlisting}

\textbf{Flusso di Esecuzione}:
\begin{enumerate}
    \item \textbf{Guard Clause}: Verifica \texttt{canSave} (previene doppi submit)
    \item \textbf{Lock State}: Imposta \texttt{isSaving = true}
    \item \textbf{Branching}: Se card temporanea, usa \texttt{onCreate}, altrimenti \texttt{onUpdate}
    \item \textbf{Success Path}: Toast di successo, exit da editing mode
    \item \textbf{Error Path}: Toast di errore, mantiene editing mode per retry
    \item \textbf{Cleanup}: \texttt{finally} block garantisce unlock di \texttt{isSaving}
\end{enumerate}

\subsubsection{handleCardClick}

\begin{lstlisting}[style=typescript, caption=Handler Click Card]
const handleCardClick = useCallback((e: React.MouseEvent) => {
    if (shouldHandleCardClick(e, isEditing) && onClick) {
        onClick(card);
    }
}, [isEditing, onClick, card]);
\end{lstlisting}

\textbf{Logica di Filtraggio}:
La funzione \texttt{shouldHandleCardClick()} previene propagazione quando:
\begin{itemize}
    \item Componente in editing mode
    \item Click su elementi interattivi (button, textarea, input)
\end{itemize}

% ============================================
% SEZIONE 4: INTERAZIONI E COMUNICAZIONE
% ============================================

\section{Interazioni e Comunicazione (I/O)}

\subsection{Comunicazione Upstream (Parent)}

Il componente comunica con il parent attraverso callbacks:

\begin{description}
    \item[\texttt{onUpdate(cardId, front, back)}] \textbf{Trigger}: Salvataggio modifiche. \textbf{Payload}: ID card e contenuti validati. \textbf{Contratto}: Promise che risolve quando persistenza completata.
    
    \item[\texttt{onCreate(front, back)}] \textbf{Trigger}: Salvataggio card temporanea. \textbf{Payload}: Solo contenuti (ID generato dal parent). \textbf{Contratto}: Promise che risolve quando card creata.
    
    \item[\texttt{onClick(card)}] \textbf{Trigger}: Click su card (view mode). \textbf{Payload}: Oggetto card completo. \textbf{Contratto}: Sincrono, nessun ritorno.
    
    \item[\texttt{onDelete(cardId)}] \textbf{Trigger}: Click pulsante delete. \textbf{Payload}: Solo ID. \textbf{Contratto}: Sincrono, gestione eliminazione delegata al parent.
    
    \item[\texttt{onCancel()}] \textbf{Trigger}: Annullamento editing. \textbf{Payload}: Nessuno. \textbf{Contratto}: Sincrono, cleanup opzionale.
    
    \item[\texttt{onShowSource(card)}] \textbf{Trigger}: Click pulsante "Show Source". \textbf{Payload}: Card completa (per accesso a \texttt{sourceMetadata}). \textbf{Contratto}: Sincrono, navigazione PDF gestita dal parent.
\end{description}

\subsection{Comunicazione Downstream (Children)}

Il componente non ha children diretti, ma utilizza sub-componenti interni:

\begin{description}
    \item[\texttt{ViewMode}] Riceve: \texttt{card}, \texttt{buttonState}, \texttt{isSourceActive}, handlers. Passa: Nessun dato, solo event handlers.
    
    \item[\texttt{EditMode}] Riceve: \texttt{tempFront}, \texttt{tempBack}, \texttt{isSaving}, \texttt{canSave}, handlers. Passa: Valori controllati e validazione state.
\end{description}

\subsection{Interazioni Esterne}

\begin{description}
    \item[\texttt{emitToast}] \textbf{Modulo}: \texttt{shared/components/toast}. \textbf{Uso}: Feedback utente per successo/errore operazioni. \textbf{Pattern}: Toast notifications non bloccanti.
    
    \item[\texttt{Card Type}] \textbf{Modulo}: \texttt{services/studyService}. \textbf{Uso}: Tipo TypeScript importato per type safety. \textbf{Contratto}: Interfaccia definita nel service layer.
\end{description}

% ============================================
% SEZIONE 5: DIAGRAMMI TECNICI
% ============================================

\section{Diagrammi Tecnici}

\subsection{Diagramma di Sequenza: Flusso Completo}

\begin{figure}[H]
\centering
\begin{tikzpicture}[
    node distance=1.5cm,
    every node/.style={align=center},
    component/.style={rectangle, draw=accent, fill=accent!10, minimum width=2.5cm, minimum height=1cm},
    action/.style={rectangle, draw=blue!50, fill=blue!5, minimum width=2cm, minimum height=0.8cm},
]

% Nodes
\node[component] (parent) {Parent\\Component};
\node[component, right=of parent] (item) {FlashcardItem};
\node[component, right=of item] (viewmode) {ViewMode};
\node[component, below=of item] (editmode) {EditMode};
\node[action, below=of parent] (user) {User\\Action};

% Initialization
\draw[->, thick] (parent) -- node[above] {\texttt{card, onUpdate}} (item);
\draw[->, thick] (item) -- node[above] {\texttt{props}} (viewmode);

% User clicks edit
\draw[->, dashed, blue] (user) -- node[left] {Click Edit} (item);
\draw[->, thick, blue] (item) -- node[right] {\texttt{setIsEditing(true)}} (editmode);

% User types
\draw[->, dashed, blue] (user) -- node[below] {Type Input} (editmode);
\draw[->, thick, blue] (editmode) -- node[left] {\texttt{setTempFront/Back}} (item);

% Validation
\draw[->, thick, green] (item) -- node[above, sloped] {\texttt{validateCardContent()}} (item);
\draw[->, thick, green] (item) -- node[right] {\texttt{canSave}} (editmode);

% User saves
\draw[->, dashed, blue] (user) -- node[below] {Click Save} (editmode);
\draw[->, thick, blue] (editmode) -- node[left] {\texttt{handleSave()}} (item);
\draw[->, thick, red] (item) -- node[above] {\texttt{onUpdate(cardId, front, back)}} (parent);

% Response
\draw[->, dashed, green] (parent) -- node[below] {\texttt{Promise.resolve()}} (item);
\draw[->, thick, green] (item) -- node[right] {\texttt{setIsEditing(false)}} (viewmode);
\draw[->, thick, green] (item) -- node[above, sloped] {\texttt{emitToast.success()}} (item);

\end{tikzpicture}
\caption{Diagramma di sequenza: Flusso completo di interazione FlashcardItem}
\label{fig:sequence}
\end{figure}

\subsection{Diagramma di Flusso: State Machine}

\begin{figure}[H]
\centering
\begin{tikzpicture}[
    node distance=2.5cm,
    state/.style={ellipse, draw=accent, fill=accent!10, minimum width=2.5cm, minimum height=1.5cm},
    transition/.style={->, thick, blue},
    decision/.style={diamond, draw=orange, fill=orange!10, minimum width=2cm, minimum height=1.5cm},
]

% States
\node[state] (view) {View\\Mode};
\node[state, right=of view] (edit) {Edit\\Mode};
\node[state, below=of view] (saving) {Saving\\State};

% Transitions
\draw[transition] (view) -- node[above] {\texttt{handleStartEdit()}} (edit);
\draw[transition] (edit) -- node[below] {\texttt{handleSave()}} (saving);
\draw[transition] (saving) -- node[left] {\texttt{Success}} (view);
\draw[transition, bend right=30] (edit) -- node[below] {\texttt{handleCancel()}} (view);
\draw[transition, bend left=30] (saving) -- node[right] {\texttt{Error}} (edit);

% Decision points
\node[decision, above=of edit] (validate) {Valid?};
\draw[transition] (edit) -- (validate);
\draw[transition] (validate) -- node[left] {No} (edit);
\draw[transition] (validate) -- node[right] {Yes} (saving);

\end{tikzpicture}
\caption{State Machine: Transizioni di stato FlashcardItem}
\label{fig:statemachine}
\end{figure}

\subsection{Diagramma Architetturale: Struttura Modulare}

\begin{figure}[H]
\centering
\begin{tikzpicture}[
    node distance=1.8cm,
    module/.style={rectangle, draw=accent, fill=accent!10, rounded corners, minimum width=3cm, minimum height=1.2cm, text centered},
    data/.style={rectangle, draw=green!50, fill=green!5, rounded corners, minimum width=2.5cm, minimum height=0.8cm},
    arrow/.style={->, thick},
]

% Main component
\node[module] (main) {FlashcardItem.tsx\\Orchestrator};

% Sub-components
\node[module, above=of main] (view) {ViewMode.tsx};
\node[module, below=of main] (edit) {EditMode.tsx};

% Supporting modules
\node[data, left=of main] (constants) {constants.ts};
\node[data, right=of main] (types) {types.ts};
\node[data, below=of constants] (utils) {utils.ts};

% Connections
\draw[arrow] (main) -- node[left] {renders} (view);
\draw[arrow] (main) -- node[left] {renders} (edit);
\draw[arrow] (constants) -- node[above] {imports} (main);
\draw[arrow] (types) -- node[above] {imports} (main);
\draw[arrow] (utils) -- node[below] {imports} (main);
\draw[arrow] (constants) -- node[above, sloped] {imports} (view);
\draw[arrow] (constants) -- node[below, sloped] {imports} (edit);

\end{tikzpicture}
\caption{Architettura modulare: Dipendenze e struttura}
\label{fig:architecture}
\end{figure}

% ============================================
% SEZIONE 6: EDGE CASES E ERROR HANDLING
% ============================================

\section{Edge Cases e Error Handling}

\subsection{Gestione Dati Mancanti}

\begin{description}
    \item[\texttt{card} null/undefined] \textbf{Prevenzione}: TypeScript garantisce prop required. \textbf{Fallback}: Componente non renderizza se card mancante (parent responsabilità).
    
    \item[\texttt{card.front/back} vuoti] \textbf{Validazione}: \texttt{validateCardContent()} verifica \texttt{length > 0}. \textbf{UI}: Pulsante Save disabilitato se campi vuoti.
    
    \item[\texttt{card.sourceMetadata} assente] \textbf{Comportamento}: Pulsante "Show Source" non renderizzato. \textbf{Logica}: \texttt{getButtonState()} ritorna \texttt{shouldShowSourceButton = false}.
\end{description}

\subsection{Gestione Errori API}

\begin{lstlisting}[style=typescript, caption=Error Handling Pattern]
try {
    await onUpdate(card.id, validation.trimmedFront, validation.trimmedBack);
    emitToast.success(...);
    setIsEditing(false);
} catch (error) {
    const errorMessage = error instanceof Error 
        ? error.message 
        : TEXT_CONTENT.toast.error.message;
    emitToast.error(errorMessage, {...});
} finally {
    setIsSaving(false);
}
\end{lstlisting}

\textbf{Strategia}:
\begin{enumerate}
    \item \textbf{Type Guard}: Verifica che error sia istanza di \texttt{Error}
    \item \textbf{Fallback Message}: Usa messaggio default se error non standard
    \item \textbf{State Preservation}: Mantiene editing mode per permettere retry
    \item \textbf{Cleanup Guaranteed}: \texttt{finally} block garantisce unlock anche su error
\end{enumerate}

\subsection{Concorrenza e Race Conditions}

\begin{description}
    \item[\texttt{isSaving} Flag] \textbf{Problema}: Doppi click su Save durante operazione asincrona. \textbf{Soluzione}: Guard clause \texttt{if (!validation.canSave) return} previene esecuzione concorrente.
    
    \item[Sincronizzazione Props] \textbf{Problema}: Card aggiornata esternamente durante editing. \textbf{Soluzione}: \texttt{useEffect} sincronizza solo se \texttt{!isEditing}, preservando modifiche in corso.
    
    \item[Card Temporanee] \textbf{Problema}: Card con ID \texttt{"temp-*"} richiede \texttt{onCreate} invece di \texttt{onUpdate}. \textbf{Soluzione}: Branching logico in \texttt{handleSave()} basato su \texttt{isTemporaryCard()}.
\end{description}

\subsection{Accessibilità e UX Edge Cases}

\begin{description}
    \item[Keyboard Navigation] \textbf{Supporto}: \texttt{autoFocus} su textarea front in EditMode. \textbf{Limiti}: Navigazione tra campi tramite Tab standard.
    
    \item[Screen Readers] \textbf{Supporto}: \texttt{aria-label} su tutti i pulsanti. \textbf{Contenuto}: Labels descrittivi per azioni (Edit, Delete, Show Source).
    
    \item[Touch Devices] \textbf{Supporto}: Pulsanti con \texttt{active:scale-95} per feedback tattile. \textbf{Area Click}: Sufficientemente grande per touch targets.
\end{description}

% ============================================
% SEZIONE 7: SUB-COMPONENTI
% ============================================

\section{Sub-Componenti}

\subsection{ViewMode Component}

\subsubsection{Responsabilità}

Il componente \texttt{ViewMode} gestisce esclusivamente la visualizzazione read-only della card.

\subsubsection{Props}

\begin{tabularx}{\textwidth}{|l|X|c|}
\hline
\textbf{Prop} & \textbf{Descrizione} & \textbf{Tipo} \\
\hline
\texttt{card} & Card da visualizzare & \texttt{Card} \\
\hline
\texttt{buttonState} & Stato dei pulsanti (calcolato) & \texttt{ButtonState} \\
\hline
\texttt{isSourceActive} & Flag highlight sorgente & \texttt{boolean} \\
\hline
\texttt{onEdit} & Handler inizio editing & \texttt{() => void} \\
\hline
\texttt{onDelete} & Handler eliminazione & \texttt{(cardId: string) => void} \\
\hline
\texttt{onShowSource} & Handler navigazione PDF & \texttt{(e: MouseEvent) => void} \\
\hline
\texttt{onCardClick} & Handler click card & \texttt{(e: MouseEvent) => void} \\
\hline
\end{tabularx}

\subsubsection{Logica di Rendering}

Il componente calcola classi CSS dinamicamente per il pulsante "Show Source":

\begin{lstlisting}[style=typescript, caption=Calcolo Classi Dinamiche]
const getShowSourceButtonClasses = (): string => {
    const base = `${BUTTON_STYLES.base.size} ...`;
    if (!buttonState.hasOnShowSource) {
        return `${base} ...disabled styles`;
    }
    if (isSourceActive) {
        return `${base} ...active styles`;
    }
    return `${base} ...default styles`;
};
\end{lstlisting}

\subsection{EditMode Component}

\subsubsection{Responsabilità}

Il componente \texttt{EditMode} gestisce l'interfaccia di modifica con validazione real-time.

\subsubsection{Props}

\begin{tabularx}{\textwidth}{|l|X|c|}
\hline
\textbf{Prop} & \textbf{Descrizione} & \textbf{Tipo} \\
\hline
\texttt{frontValue} & Valore temporaneo front & \texttt{string} \\
\hline
\texttt{backValue} & Valore temporaneo back & \texttt{string} \\
\hline
\texttt{isSaving} & Flag operazione in corso & \texttt{boolean} \\
\hline
\texttt{canSave} & Flag validazione passata & \texttt{boolean} \\
\hline
\texttt{onFrontChange} & Handler cambio front & \texttt{(value: string) => void} \\
\hline
\texttt{onBackChange} & Handler cambio back & \texttt{(value: string) => void} \\
\hline
\texttt{onSave} & Handler salvataggio & \texttt{() => void} \\
\hline
\texttt{onCancel} & Handler annullamento & \texttt{() => void} \\
\hline
\end{tabularx}

\subsubsection{Loading State}

Durante salvataggio, il componente mostra spinner animato:

\begin{lstlisting}[style=typescript, caption=Loading Spinner]
{saving ? (
    <>
        <motion.div
            animate={{ rotate: 360 }}
            transition={ANIMATION_CONFIG.loadingSpinner}
            className="w-4 h-4 border-2 border-white/30 border-t-white rounded-full"
        />
        <span>{TEXT_CONTENT.buttons.saving}</span>
    </>
) : (
    <>
        <FiCheck className={ICON_SIZES.small} />
        <span>{TEXT_CONTENT.buttons.save}</span>
    </>
)}
\end{lstlisting}

% ============================================
% SEZIONE 8: PERFORMANCE E OTTIMIZZAZIONI
% ============================================

\section{Performance e Ottimizzazioni}

\subsection{React.memo}

Il componente è wrappato in \texttt{React.memo} per prevenire re-render inutili:

\begin{lstlisting}[style=typescript, caption=Memoization]
export const FlashcardItem: React.FC<FlashcardItemProps> = memo(({ ... }) => {
    // Component implementation
});
\end{lstlisting}

\textbf{Benefici}: Re-render solo quando props cambiano (shallow comparison).

\subsection{useCallback Optimization}

Tutti gli event handlers utilizzano \texttt{useCallback} con dipendenze ottimizzate:

\begin{itemize}
    \item \texttt{handleStartEdit}: Dipende da \texttt{card.front}, \texttt{card.back}
    \item \texttt{handleSave}: Dipende da \texttt{card.id}, \texttt{validation}, callbacks
    \item \texttt{handleCardClick}: Dipende da \texttt{isEditing}, \texttt{onClick}, \texttt{card}
\end{itemize}

\textbf{Impatto}: Previene ricreazione funzioni ad ogni render, riducendo re-render di children.

\subsection{Computed Values}

Valori computati sono calcolati una volta per render:

\begin{lstlisting}[style=typescript, caption=Computed Values]
const validation = validateCardContent(...);
const buttonState = getButtonState(card, onShowSource);
\end{lstlisting}

\textbf{Pattern}: Funzioni pure senza side effects, facilmente testabili.

% ============================================
% CONCLUSIONI
% ============================================

\section{Conclusioni}

Il componente \texttt{FlashcardItem} rappresenta un'implementazione robusta e production-ready di un'interfaccia dual-mode per la gestione di flashcard. L'architettura modulare, la separazione delle responsabilità e l'uso di pattern React ottimizzati garantiscono manutenibilità, performance e type safety.

\textbf{Punti di Forza}:
\begin{itemize}
    \item Architettura modulare con separazione concerns
    \item Type safety completo con TypeScript
    \item Gestione errori robusta con fallback
    \item Performance ottimizzate con memoization
    \item Accessibilità e UX curate
\end{itemize}

\textbf{Aree di Miglioramento Future}:
\begin{itemize}
    \item Aggiunta unit tests per funzioni pure
    \item Implementazione undo/redo per editing
    \item Supporto markdown nei campi testo
    \item Ottimizzazione animazioni per dispositivi low-end
\end{itemize}

\end{document}
